\section*{第 \ref{ch:b1:nucleic-acids} 章——核苷酸与核酸}

\begin{solution}{exo:b1:nucleic-acids:1}
含氮\omterm{def:b1:nucleic-acids:nucleotide}{碱基}、戊糖、一至三个磷酸基团。\omterm{def:b1:nucleic-acids:chain}{RNA} 中为核糖
（带 $2'$ 羟基）和尿\omterm{def:b1:nucleic-acids:nucleotide}{嘧啶}；\omterm{def:b1:nucleic-acids:chain}{DNA} 中为 $2'$-脱氧核糖和胸腺\omterm{def:b1:nucleic-acids:nucleotide}{嘧啶}。
\end{solution}

\begin{solution}{exo:b1:nucleic-acids:2}
$5'$-TAGGATCC-$3'$。配对：G--C、G--C、A--T、T--A、C--G、C--G、T--A、
A--T：四对 G--C（12 个\omterm{def:b1:water-small-molecules:water}{氢键}）和四对 A--T（8 个），共 20 个\omterm{def:b1:water-small-molecules:water}{氢键}。
\end{solution}

\begin{solution}{exo:b1:nucleic-acids:3}
T 占 \qty{22}{\%}；G 与 C 共占其余的 \qty{56}{\%}，各占
\qty{28}{\%}；$G + C = \qty{56}{\%}$。
\end{solution}

\begin{solution}{exo:b1:nucleic-acids:4}
在两条曲线之间，按每百分点 \qty{0.4}{\celsius} 计，约为
\qty{85}{\celsius}。吸光度升高 \qty{40}{\%}。
\end{solution}

\begin{solution}{exo:b1:nucleic-acids:5}
$\num{4.6e6}\times 0.34\,\mathrm{nm} = \qty{1.56}{mm}$; \num{460000}
圈；长度压缩 $1560/2 = 780$ 倍（染色体折叠为
环和超螺旋）。
\end{solution}

\begin{solution}{exo:b1:nucleic-acids:6}
\omterm{def:b1:nucleic-acids:nucleotide}{碱基}以边缘相对的方式配对，各\omterm{def:b1:nucleic-acids:nucleotide}{碱基}所连的糖
只有两条骨架方向相反时，才会位于\omterm{thm:b1:nucleic-acids:helix}{碱基对}的同一侧；
若两条链平行，两个糖将位于
\omterm{thm:b1:nucleic-acids:helix}{碱基对}的相对两侧，\omterm{thm:b1:nucleic-acids:helix}{碱基对}的宽度以及骨架的位置
都会随序列而变。反向平行的两条链与\omterm{def:b1:nucleic-acids:nucleotide}{嘌呤}—\omterm{def:b1:nucleic-acids:nucleotide}{嘧啶}配对
给出恒定的 \qty{2}{nm}
宽度，这正是纤维衍射图样所需。
\end{solution}

\begin{solution}{exo:b1:nucleic-acids:7}
$0.35\times 50 = \qty{17.5}{\micro g/mL}$. $A_{260}/A_{280} = 1.75$:
基本为纯 \omterm{def:b1:nucleic-acids:chain}{DNA}（1.8），可能含有微量蛋白质。
在 \qty{1}{mL} 中：$17.5\times 10^{-6}/(650\times 1.66\times 10^{-24}) =
\num{1.6e16}$ 个\omterm{thm:b1:nucleic-acids:helix}{碱基对}。
\end{solution}

\begin{solution}{exo:b1:nucleic-acids:8}
堆叠的\omterm{def:b1:nucleic-acids:nucleotide}{碱基}相互遮挡，吸光较少；分开两链
使每个\omterm{def:b1:nucleic-acids:nucleotide}{碱基}暴露于光线下，吸光度升高
\qty{40}{\%}。纯 \omterm{def:b1:nucleic-acids:chain}{DNA} 在其单一的
$T_m$ 附近几度之内协同熔化；混合物含有许多 $G + C$ 含量各异的 \omterm{def:b1:nucleic-acids:chain}{DNA}，
每种都有各自的 $T_m$，因此其曲线之和
分布在数十度的范围内。
\end{solution}

\begin{solution}{exo:b1:nucleic-acids:9}
略低于 \qty{62}{\celsius} 且高于 \qty{57}{\celsius}，例如
\qty{60}{\celsius}：完全配对的\omterm{def:b1:nucleic-acids:melting}{杂交}分子保持稳定，错配的
则熔化。在 \qty{45}{\celsius} 时，含至多三个错配的\omterm{def:b1:nucleic-acids:melting}{杂交}分子
仍然稳定，探针也会使相关序列显示信号。
\end{solution}

\begin{solution}{exo:b1:nucleic-acids:10}
这些规则源于两条互补链之间的\omterm{def:b1:nucleic-acids:nucleotide}{碱基}配对。
\omterm{def:b1:nucleic-acids:chain}{RNA} 为单链，没有互补链来约束其组成，因此
一般情况下 $A \neq U$；所讨论的小病毒具有
单链 \omterm{def:b1:nucleic-acids:chain}{DNA} 基因组，同样不必遵守上述规则——
它们不遵守这些规则，曾是最早提示其为
单链的线索之一。
\end{solution}

\begin{solution}{exo:b1:nucleic-acids:11}
单个\omterm{thm:b1:nucleic-acids:helix}{碱基对}的结合能仅约每摩尔几千焦，小于热运动的
能量，且在微秒量级内形成与断裂。在螺旋中，\omterm{thm:b1:nucleic-acids:helix}{碱基对}
协同形成：每一对堆叠在前一对之上，最初几对
引发其余部分的形成，因此能量叠加，所有\omterm{thm:b1:nucleic-acids:helix}{碱基对}同时断裂的
概率可忽略不计——即协同性。堆叠作用的贡献与
\omterm{def:b1:water-small-molecules:water}{氢键}相当。对探针而言，这意味着其稳定性
随长度增加，并且每出现一个错配都会急剧下降，因为错配打断了
协同的连续区段：专一性是整段序列的性质。
\end{solution}

\begin{solution}{exo:b1:nucleic-acids:12}
\omterm{thm:b1:nucleic-acids:helix}{双螺旋}结构立即解释了复制：每条链决定了另一条链的
序列，因此将两条链分开，并在各自链上配对新的\omterm{def:b1:nucleic-acids:nucleotide}{核苷酸}，便得到两个
相同的分子；突变即某个拷贝中固定下来的错误\omterm{def:b1:nucleic-acids:nucleotide}{碱基}，
此后被忠实地复制。它本身并不能说明
四种\omterm{def:b1:nucleic-acids:nucleotide}{碱基}的序列如何决定蛋白质——这花了十年时间
（遗传密码、mRNA、tRNA、核糖体）才得以弄清。该模型被接受，是因为它
与所有测量结果相符（\omterm{prop:b1:nucleic-acids:chargaff}{查加夫}规则、X 射线测得的尺寸、密度），
没有任何剩余，而且它对遗传的解释非常
简洁，以至于另一种可能——纯属巧合——显得难以置信；
随后梅塞尔森与斯塔尔证实了它所预言的复制方式。
\end{solution}

\begin{solution}{pb:b1:nucleic-acids:1}
\textbf{1.} $\num{48502}\times 0.34 = \qty{16490}{nm} = \qty{16.5}{\micro m}$.
\textbf{2.} $\num{48502}/10 = 4850$ 圈。
\textbf{3.} $\num{48502}\times 650 = \qty{3.15e7}{Da}$.
\textbf{4.} $\num{3.15e7}\times 1.66\times 10^{-24} = \qty{5.2e-17}{g}$.
\textbf{5.} 每对两个磷酸基团，共 \num{97004} 个，净电荷 $-97\,004\,e$。
\textbf{6.} $G + C = 0.499\times\num{97004} = \num{48405}$: $G = C =
\num{24203}$; $A = T = \num{24299}$.
\textbf{7.} $3\times\num{24203} + 2\times\num{24299} = \num{121207}$.
\textbf{8.} $81.5 + 16.6\times(-1) + 0.41\times 49.9 - 500/48502 =
81.5 - 16.6 + 20.5 - 0.01 = \qty{85.4}{\celsius}$.
\textbf{9.} $81.5 - 33.2 + 20.5 = \qty{68.8}{\celsius}$。阳离子屏蔽
两条带负电的骨架之间的斥力；盐较少时，
两链排斥更强，在更低温度下即分开。
\textbf{10.} \qty{35}{\%}: $81.5 - 16.6 + 14.4 - 0.5 =
\qty{78.8}{\celsius}$; \qty{62}{\%}: $81.5 - 16.6 + 25.4 - 0.5 =
\qty{89.8}{\celsius}$。整个分子分步熔化：富含 A--T 的区域先
解开，富含 G--C 的区域最后解开，因此曲线呈宽的
上升并带肩，而非单一的陡峭台阶。
\textbf{11.} $0.50\times 50 = \qty{25}{\micro g/mL}$；熔化后
$A_{260} = 0.70$.
\textbf{12.} $25\times 10^{-6}/5.2\times 10^{-17} = \num{4.8e11}$ 个分子。
\textbf{13.} $81.5 - 16.6 + 0.41\times 50 - 500/12 = 81.5 - 16.6 + 20.5
- 41.7 = \qty{43.7}{\celsius}$（该公式用于如此短的
片段仅为粗略估计）。在 \qty{37}{\celsius} 下两端可配对，但会张开；
而在\omterm{def:b1:cell-unit-of-life:cell}{细胞}中，由连接酶共价封闭两处切口，使
环状分子不再依赖这 12 对\omterm{def:b1:nucleic-acids:nucleotide}{碱基}。
\textbf{14.} $\pi\times(1.0)^2\times\num{16490} = \qty{5.2e4}{nm^3}$.
\textbf{15.} $\frac{4}{3}\pi\times 27.5^3 = \qty{8.7e4}{nm^3}$.
\textbf{16.} \qty{60}{\%}：约为最紧密可能\omterm{prop:b1:nucleic-acids:stability}{堆积}的三分之二
（圆柱体）的\omterm{prop:b1:nucleic-acids:stability}{堆积}——\omterm{def:b1:nucleic-acids:chain}{DNA} 以近似晶体的卷曲方式缠绕。
\textbf{17.} 抗衡离子（多胺、$\mathrm{Mg^{2+}}$）必须
中和大部分电荷，否则无法克服斥力；
残余的斥力，以及紧密卷曲中储存的弯曲能，
向外推动外壳；尾巴打开时，这一压力
将基因组的前一部分推入细菌中。
\textbf{18.} \qty{55}{nm} 衣壳中的环半径为
\qtyrange{10}{25}{nm}，低于 \omterm{def:b1:nucleic-acids:chain}{DNA} 可以自由弯曲的
\qty{50}{nm} 这一临界值：弯曲需耗能，由消耗 \omterm{def:b1:water-small-molecules:atp}{ATP} 的包装马达提供，
并以第 17 题的压力形式储存。
\textbf{19.} $5'$-AGGTCGCCGCCC-$3'$.
\textbf{20.} \num{1200} 个\omterm{def:b1:nucleic-acids:nucleotide}{核苷酸}；$1200\times 330 = \qty{4.0e5}{Da}$。
\textbf{21.} $75\times 600 = \num{45000}$ 对，\qty{93}{\%} 为编码序列
——基因组紧凑；人类基因组在
基因密度上约低 60 倍。
\textbf{22.} 两个复制叉，各在 \qty{1200}{s} 内复制 \num{24251} 对：
每叉每秒 \qty{20}{} 对（宿主聚合酶的能力可达
\qty{1000}{}；耗时由整个周期决定，而非由
聚合酶决定）。
\textbf{23.} $\num{121207}/\num{48502} = 2.5$ 个\omterm{def:b1:water-small-molecules:water}{氢键}每对；两个复制叉各以每秒 20 对推进，每秒断开 100 个\omterm{def:b1:water-small-molecules:water}{氢键}，其后方同样数目的\omterm{def:b1:water-small-molecules:water}{氢键}随之重新形成。
\textbf{24.} \num{48502} 对中仅一对发生改变时，$G + C$ 变化
\qty{0.002}{\%}，$T_m$ 变化 \qty{0.001}{\celsius}：无法通过
整体熔化分子来检测；只有跨越该位点的短探针
才能显示它。
\textbf{25.} $T_m = \qty{85}{\celsius}$（在 \qty{0.1}{mol/L} 钠盐中）；
轮廓长度 \qty{16.5}{\micro m}，为衣壳
直径的 300 倍，装入其中时占其体积的 \qty{60}{\%}。
\end{solution}
